\documentclass[11pt,reqno]{amsart}

% ================================================================
% Packages
% ================================================================
\usepackage{amsmath,amssymb,amsthm}
\usepackage{mathrsfs}
\usepackage{enumerate}
\usepackage[shortlabels]{enumitem}
\usepackage{booktabs}
\usepackage{array}
\usepackage{microtype}
\usepackage[dvipsnames]{xcolor}
\usepackage[colorlinks=true,
  linkcolor=blue!60!black,
  citecolor=green!40!black,
  urlcolor=blue!60!black]{hyperref}

% ================================================================
% Theorem environments
% ================================================================
\newtheorem{theorem}{Theorem}[section]
\newtheorem{proposition}[theorem]{Proposition}
\newtheorem{lemma}[theorem]{Lemma}
\newtheorem{corollary}[theorem]{Corollary}
\newtheorem{conjecture}[theorem]{Conjecture}
\theoremstyle{definition}
\newtheorem{definition}[theorem]{Definition}
\newtheorem{construction}[theorem]{Construction}
\newtheorem{example}[theorem]{Example}
\newtheorem{warning}[theorem]{Warning}
\theoremstyle{remark}
\newtheorem{remark}[theorem]{Remark}

% ================================================================
% Macros
% ================================================================
\newcommand{\cA}{\mathcal{A}}
\newcommand{\cB}{\mathcal{B}}
\newcommand{\cC}{\mathcal{C}}
\newcommand{\cD}{\mathcal{D}}
\newcommand{\cF}{\mathcal{F}}
\newcommand{\cH}{\mathcal{H}}
\newcommand{\cM}{\mathcal{M}}
\newcommand{\cO}{\mathcal{O}}
\newcommand{\cL}{\mathcal{L}}
\newcommand{\cP}{\mathcal{P}}
\newcommand{\cQ}{\mathcal{Q}}
\newcommand{\cW}{\mathcal{W}}
\newcommand{\cZ}{\mathcal{Z}}

\newcommand{\Ran}{\mathrm{Ran}}
\newcommand{\MC}{\mathrm{MC}}
\newcommand{\HH}{\mathrm{HH}}
\newcommand{\ChirHoch}{\mathrm{ChirHoch}}
\newcommand{\Ainf}{A_{\infty}}

\newcommand{\Eone}{\mathsf{E}_{1}}
\newcommand{\Etwo}{\mathsf{E}_{2}}
\newcommand{\Ethree}{\mathsf{E}_{3}}
\newcommand{\En}{\mathsf{E}_{n}}
\newcommand{\Einf}{\mathsf{E}_{\infty}}

\newcommand{\Sym}{\mathrm{Sym}}
\newcommand{\Ass}{\mathrm{Ass}}
\newcommand{\Com}{\mathrm{Com}}
\newcommand{\Lie}{\mathrm{Lie}}
\newcommand{\Vir}{\mathrm{Vir}}
\newcommand{\Rep}{\mathrm{Rep}}
\newcommand{\Sh}{\mathrm{Sh}}
\newcommand{\DS}{\mathrm{DS}}
\newcommand{\SCchtop}{\mathrm{SC}^{\mathrm{ch,top}}}
\newcommand{\Convstr}{\mathrm{Conv}_{\mathrm{str}}}
\newcommand{\barB}{\bar{B}}
\newcommand{\FM}{\overline{C}}
\newcommand{\hol}{\mathrm{hol}}

\newcommand{\fg}{\mathfrak{g}}
\newcommand{\fh}{\mathfrak{h}}
\newcommand{\fsl}{\mathfrak{sl}}
\newcommand{\gmod}{\mathfrak{g}^{\mathrm{mod}}_{\cA}}

\newcommand{\bC}{\mathbb{C}}
\newcommand{\bR}{\mathbb{R}}
\newcommand{\bZ}{\mathbb{Z}}
\newcommand{\bQ}{\mathbb{Q}}

\DeclareMathOperator{\Hom}{Hom}
\DeclareMathOperator{\RHom}{RHom}
\DeclareMathOperator{\Res}{Res}
\DeclareMathOperator{\Aut}{Aut}
\DeclareMathOperator{\End}{End}
\DeclareMathOperator{\Conf}{Conf}
\DeclareMathOperator{\Spec}{Spec}
\DeclareMathOperator{\av}{av}
\DeclareMathOperator{\rk}{rk}
\DeclareMathOperator{\tr}{tr}
\DeclareMathOperator{\coll}{coll}
\DeclareMathOperator{\gr}{gr}
\DeclareMathOperator{\cofib}{cofib}

\numberwithin{equation}{section}

% ================================================================
\begin{document}

\title{The holographic modular Koszul datum}

\author{Raeez Lorgat}
\date{}

\begin{abstract}
Every 3d holomorphic-topological gauge theory $T$ on
$\bR_+ \times X$ is controlled by a six-component package
$\cH(T) = (\cA, \cA^!, \cC, r(z), \Theta_{\cA},
\nabla^{\hol})$: boundary chiral algebra, Koszul dual,
line-operator category, spectral $r$-matrix, universal
Maurer--Cartan element, and modular shadow connection. We define
each component from the bar complex
$\barB(\cA) = T^c(s^{-1}\bar{\cA})$ and its single
structural identity $D_{\cA}^2 = 0$, and compute the datum in
full for three families: Heisenberg $\cH_k$ (abelian
Chern--Simons, class $\mathbf{G}$), affine
$\widehat{\fsl}_2{}_k$ (non-abelian Chern--Simons, class
$\mathbf{L}$), and Virasoro $\Vir_c$ (3d gravity, class
$\mathbf{M}$). The bulk--boundary--line triangle is recovered
from the datum: the bulk is the chiral derived centre
$\cZ^{\mathrm{der}}_{\mathrm{ch}}(\cA)$, the boundary is
$\cA$ itself, and the line category is $\cA^!\text{-mod}$. The
$\SCchtop$ structure emerges on the pair
$(\cZ^{\mathrm{der}}_{\mathrm{ch}}(\cA), \cA)$ of derived centre
and boundary algebra, not on the bar complex. The central
organising principle is the $\SCchtop$ \emph{heptagon}: seven edges
(operadic / Koszul / factorization / BV--BRST / convolution / Drinfeld
centre / derived-AG) identify the datum across presentations, with
six classical faces and a seventh derived-algebraic-geometry face
inscribed via PTVV on $\mathrm{Map}(X \times \bR_{\geq 0},
B\,\mathrm{SC\text{-}Alg})$. We construct the
anomaly completion: the transgression algebra $B_\Theta$ with
secondary anomaly $u = \eta^2 = \kappa(\cA^!) \cdot \omega_g$,
whose vanishing at the critical central charge is the
Clifford/exterior bifurcation governing the genus tower.
\end{abstract}

\begin{remark}[Wave-14 reconstitution anchor]\label{rem:wave14-anchor-HD}
This standalone is the inscription of the Platonic holographic-datum
heptagon. The closure wave (2026-04-16) added two load-bearing edges
to the earlier five-face star. Face~(6) is the Drinfeld-centre
identification
$Z(\mathsf{Rep}_{\mathrm{fact}}(\cA)) \simeq
\mathsf{Rep}_{\mathrm{fact}}(\cZ^{\mathrm{der}}_{\mathrm{ch}}(\cA))^{\Etwo}$
via categorified bar-cobar with half-braiding (cross-ref
\textup{\textsf{thm:universal-holography-master}}, Vol~II
\verb|chapters/theory/sc_chtop_heptagon.tex|). Face~(7) is the
derived-algebraic-geometry edge via PTVV 1-shifted symplectic on
$\mathrm{Map}(X \times \bR_{\geq 0}, B\,\mathrm{SC\text{-}Alg})$.
Every face is genuinely independent (no tautological routing); the
seven classical identifications (operadic, Koszul, factorization,
BV--BRST, convolution, Drinfeld-centre, derived-AG) all factor through
the chiral Hochschild cochain complex as hub. The topologisation from
$\SCchtop$ to $\Ethree^{\mathrm{top}}$ is PROVED chain-level on the
original complex for classes $G$, $L$, $C$, and the critical level
(Vol~II \verb|chapters/theory/sc_chtop_heptagon.tex|,
\textup{\textsf{thm:topologization-tower}}); for class~$M$ it is
weight-completed (correct scope, direct-sum chain-level genuinely
fails).
\end{remark}

\maketitle

\tableofcontents


% ====================================================================
\section{Introduction}\label{sec:intro}
% ====================================================================

The problem is to organise the algebraic data of a 3d
holomorphic-topological (HT) quantum field theory into a single
invariant from which every physical observable reconstructs. A 3d HT
theory on the half-space $\bR_+ \times X$, where $X$ is a smooth
algebraic curve, produces three observation algebras: a boundary
chiral algebra $\cA$ at $\{0\} \times X$, a bulk algebra in the
interior, and a category of topological line operators along $\bR_+$.
The algebraic engine relating these three is Koszul duality in the
homotopical modular chiral realm on algebraic curves: the bar
construction on $\cA$ produces a coalgebra from which the Koszul dual,
the $r$-matrix, and the genus tower are extracted. The governing
identity is $D_{\cA}^2 = 0$; the governing equation is the
Maurer--Cartan equation $D\Theta + \frac{1}{2}[\Theta,\Theta] = 0$.

The holographic modular Koszul datum is
\begin{equation}\label{eq:datum-intro}
\cH(T) \;=\;
(\cA,\; \cA^!,\; \cC,\; r(z),\; \Theta_{\cA},\; \nabla^{\hol}).
\end{equation}
One six-component package, extracted from a single Maurer--Cartan
element by a sequence of projections. The components are:
\begin{enumerate}[label=(\roman*),nosep]
\item the boundary chiral algebra $\cA$;
\item the chiral Koszul dual $\cA^! = (H^*(\barB(\cA)))^\vee$;
\item the line-operator category $\cC \simeq \cA^!\text{-mod}$;
\item the classical $r$-matrix
  $r(z) = \Res^{\coll}_{0,2}(\Theta_{\cA})$;
\item the universal Maurer--Cartan element $\Theta_{\cA} := D_{\cA} - d_0$;
\item the modular shadow connection
  $\nabla^{\hol}_{g,n} = d - \Sh_{g,n}(\Theta_{\cA})$.
\end{enumerate}
The bar complex $\barB(\cA) = T^c(s^{-1}\bar{\cA})$ is an $\Eone$-chiral
coassociative coalgebra: it carries a differential (from the chiral
product) and a deconcatenation coproduct (from the tensor coalgebra
structure). The $\SCchtop$ structure does not live on $\barB(\cA)$;
it emerges on the chiral derived centre
$\cZ^{\mathrm{der}}_{\mathrm{ch}}(\cA)
= C^\bullet_{\mathrm{ch}}(\cA,\cA)$, where the pair
$(\cZ^{\mathrm{der}}_{\mathrm{ch}}(\cA), \cA)$ is the SC datum
with bulk acting on boundary.

Three families illustrate the full range of the datum.

\smallskip
\noindent\emph{Heisenberg $\cH_k$.}\enspace
Abelian Chern--Simons at level $k$. Shadow class $\mathbf{G}$, shadow
depth $2$. The $r$-matrix is $r(z) = k/z$ (vanishing at $k=0$);
the line category is $\mathsf{Vect}$; the shadow connection is
trivial. This is the base case: all six components are scalar.

\smallskip
\noindent\emph{Affine $\widehat{\fsl}_2{}_k$.}\enspace
Non-abelian Chern--Simons at level $k$. Shadow class $\mathbf{L}$,
shadow depth $3$. The $r$-matrix is Yang's
$r(z) = k\,\Omega/z$ (trace-form; vanishing at $k=0$). The line
category is $\Rep(U_q(\fsl_2))$ at $q = e^{i\pi/(k+2)}$, recovered
from the KZ monodromy. The shadow connection is the KZ connection.
Every component is matrix-valued.

\smallskip
\noindent\emph{Virasoro $\Vir_c$.}\enspace
3d gravity. Shadow class $\mathbf{M}$, infinite depth. The $r$-matrix
is $r_c(z) = (c/2)/z^3 + 2T/z$ (cubic and simple poles); the quartic
contact invariant $Q^{\mathrm{contact}} = 10/[c(5c+22)] \neq 0$. The
Koszul dual is $\Vir_{26-c}$, self-dual at $c = 13$. The genus tower
is infinite; gravity does not truncate.

\smallskip

The bulk--boundary--line (BBL) triangle is not an ad hoc
diagram connecting three independent objects. It is a single object
(the bar complex with its $\Eone$-chiral coassociative structure)
viewed from three angles. The three edges are: boundary$\to$bulk
(the bar map followed by the trace to the ambient complex),
bulk$\to$lines (the complementarity cofibre), and boundary$\to$lines
(Koszul duality). Commutativity of the triangle is Theorem~A.

The anomaly completion adjoins to the Koszul dual $\cA^!$ a degree-$1$
generator $\eta$ with Clifford (not exterior) relation
$\eta^2 = \kappa(\cA^!) \cdot \omega_g$. The resulting transgression
algebra $B_\Theta$ has a secondary anomaly
$u = \eta^2 = \kappa(\cA^!) \cdot \omega_g$ that controls the genus
tower: when $u \neq 0$ the Clifford algebra is a matrix algebra and
genera are Morita-equivalent; when $u = 0$ the handle operators become
exterior generators and the theory sees genuine surface topology.

\smallskip
\noindent\textbf{Conventions.}\enspace
Grading is cohomological ($|d| = +1$). Desuspension lowers degree:
$|s^{-1}v| = |v| - 1$. The bar propagator is $d\log E(z,w)$,
weight $1$ in both variables regardless of field weight. OPE-mode
notation: $a_{(n)}b$. The augmentation ideal is
$\bar{\cA} = \ker(\varepsilon)$; the bar complex uses $\bar{\cA}$,
not $\cA$.


% ====================================================================
\section{The datum}\label{sec:datum}
% ====================================================================

Let $\cA$ be a chirally Koszul vertex algebra on a smooth projective
curve $X$.

\begin{definition}[The holographic modular Koszul datum]
\label{def:holographic-datum}
The \emph{holographic modular Koszul datum} of $\cA$ is the
six-component package
\begin{equation}\label{eq:datum}
\cH(\cA) \;=\;
(\cA,\; \cA^!,\; \cC,\; r(z),\; \Theta_{\cA},\; \nabla^{\hol}),
\end{equation}
with components defined in
Definitions~\ref{def:boundary}--\ref{def:shadow-connection}
below.
\end{definition}


\subsection{Boundary algebra}\label{subsec:boundary}

\begin{definition}[Boundary algebra]\label{def:boundary}
The \emph{boundary algebra} is the chiral algebra $\cA$ on $X$,
restricted from the holomorphic boundary
$\{0\} \times X \subset \bR_+ \times X$ of the 3d HT system.
\end{definition}

The boundary carries the full OPE data and is the input to the bar
construction. In the 3d HT picture, $\cA$ is the boundary condition
at the $t = 0$ wall: the holomorphic direction produces the OPE, and
the topological direction produces the ordering.


\subsection{Bar complex}\label{subsec:bar}

The ordered bar complex is
\begin{equation}\label{eq:bar}
\barB^{\mathrm{ord}}(\cA) \;=\; T^c(s^{-1}\bar{\cA}),
\end{equation}
where $\bar{\cA} = \ker(\varepsilon)$ is the augmentation ideal and
$s^{-1}$ denotes desuspension ($|s^{-1}v| = |v| - 1$). This is an
$\Eone$-chiral coassociative coalgebra: the differential $d_{\barB}$
extracts OPE residues at collision divisors of the
Fulton--MacPherson space $\FM_n(\bC)$, and the deconcatenation
coproduct
\begin{equation}\label{eq:deconc}
\Delta(a_1 \otimes \cdots \otimes a_n) \;=\;
\sum_{p=0}^{n}
(a_1 \otimes \cdots \otimes a_p) \otimes
(a_{p+1} \otimes \cdots \otimes a_n)
\end{equation}
splits an ordered tensor sequence at a cut point.

The symmetric bar $\barB^{\Sigma}(\cA)$ is the
$\Sigma_n$-coinvariant image under the averaging map
$\av \colon \barB^{\mathrm{ord}} \to \barB^{\Sigma}$. The
descent from ordered to symmetric is $R$-matrix-twisted:
$\barB^{\Sigma}_n = (\barB^{\mathrm{ord}}_n)^{R\text{-}\Sigma_n}$.
The ordered bar is the primitive object; the symmetric bar is
its coinvariant shadow.

\begin{warning}[Bar complex is not an SC-coalgebra]
\label{warn:bar-not-sc}
The bar complex $\barB(\cA)$ is an $\Eone$-chiral coassociative
coalgebra. It does \emph{not} carry an $\SCchtop$ structure. The
$\SCchtop$ structure emerges on the chiral derived centre
$\cZ^{\mathrm{der}}_{\mathrm{ch}}(\cA) =
C^\bullet_{\mathrm{ch}}(\cA,\cA)$, computed using the bar complex
as a resolution. The pair
$(\cZ^{\mathrm{der}}_{\mathrm{ch}}(\cA), \cA)$ is the SC datum.
\end{warning}


\subsection{Koszul dual}\label{subsec:koszul-dual}

\begin{definition}[Chiral Koszul dual]\label{def:koszul-dual}
The \emph{chiral Koszul dual} of $\cA$ is
\begin{equation}\label{eq:kd}
\cA^! \;=\; \bigl(H^*(\barB(\cA))\bigr)^\vee,
\end{equation}
the graded linear dual of the bar cohomology coalgebra
$\cA^{\mathrm{i}} = H^*(\barB(\cA))$.
\end{definition}

The passage from $\cA$ to $\cA^!$ involves three steps: bar
construction ($\cA \mapsto \barB(\cA)$, coalgebra), cohomology
($\barB(\cA) \mapsto \cA^{\mathrm{i}}$, dual coalgebra), and
Verdier duality ($\cA^{\mathrm{i}} \mapsto \cA^! =
(\cA^{\mathrm{i}})^\vee$, dual algebra). These are three distinct
functors producing four distinct objects:
\begin{equation}\label{eq:four-objects}
\cA \quad\longmapsto\quad
\barB(\cA) \quad\longmapsto\quad
\cA^{\mathrm{i}} = H^*(\barB(\cA)) \quad\longmapsto\quad
\cA^! = (\cA^{\mathrm{i}})^\vee.
\end{equation}
The cobar construction $\Omega(\barB(\cA)) \simeq \cA$ is
\emph{inversion}, not Koszul duality.


\subsection{Line-operator category}\label{subsec:lines}

\begin{definition}[Line-operator category]\label{def:lines}
The \emph{line-operator category} is
\begin{equation}\label{eq:lines}
\cC \;\simeq\; \cA^!\text{-}\mathsf{mod}.
\end{equation}
\end{definition}

Objects are topological defect lines along the $\bR_+$-direction;
morphisms are junction operators. The braiding on $\cC$ comes from
crossing in the holomorphic plane: two line operators separated in
$\bC_z$ can be exchanged by moving one around the other. On the
Koszul locus, the $r$-matrix $r(z)$ determines the braiding
explicitly.


\subsection{Spectral $r$-matrix}\label{subsec:r-matrix}

\begin{definition}[Classical $r$-matrix]\label{def:r-matrix}
The \emph{spectral $r$-matrix} is the binary genus-$0$ collision
residue
\begin{equation}\label{eq:r-matrix}
r(z) \;=\; \Res^{\coll}_{0,2}(\Theta_{\cA})
\;\in\;
\cA^! \,\widehat{\otimes}\, \cA^![\![z^{-1}, z]\!].
\end{equation}
\end{definition}

The $r$-matrix satisfies the classical Yang--Baxter equation (CYBE),
which follows from the Arnold relation on $\overline{M}_{0,4}$: the
three boundary divisors of the genus-$0$ $4$-pointed moduli space
correspond to channels $(12|34)$, $(13|24)$, $(14|23)$, and evaluating
the Arnold relation on
$\Omega \otimes \Omega \in (\cA^!)^{\otimes 3}$ produces the
three terms of $\sum [r_{ij}, r_{ik}] = 0$.

The $r$-matrix lives in the $\Eone$ convolution algebra
$\mathfrak{g}^{\Eone}_{\cA}$; the modular characteristic
$\kappa(\cA)$ is its $\Sigma_2$-coinvariant projection. For abelian
algebras, $\av(r(z)) = \kappa$ directly. For non-abelian affine
Kac--Moody algebras, the full modular characteristic includes the
Sugawara shift:
\begin{equation}\label{eq:av-sugawara}
\kappa(V_k(\fg))
\;=\;
\av(r(z)) + \frac{\dim(\fg)}{2}
\;=\;
\frac{\dim(\fg)(k + h^\vee)}{2h^\vee},
\end{equation}
where the $\dim(\fg)/2$ term is the simple-pole self-contraction
through the adjoint Casimir eigenvalue $2h^\vee$.


\subsection{Universal Maurer--Cartan element}\label{subsec:mc}

\begin{definition}[Universal MC element]\label{def:mc}
The \emph{universal Maurer--Cartan element} is
\begin{equation}\label{eq:mc}
\Theta_{\cA} \;:=\; D_{\cA} - d_0
\;\in\;
\MC(\gmod),
\end{equation}
where $D_{\cA}$ is the full bar differential and $d_0$ is its
linear (genus-$0$, degree-$1$) part.
\end{definition}

The element $\Theta_{\cA}$ is always MC because $D_{\cA}^2 = 0$:
\begin{equation}\label{eq:mc-eq}
D\Theta_{\cA} + \tfrac{1}{2}[\Theta_{\cA}, \Theta_{\cA}] = 0.
\end{equation}
This is the bar-intrinsic construction: $\Theta_{\cA}$ exists for
every chirally Koszul algebra, at all degrees, as an inverse-limit
$\Theta_{\cA} = \varprojlim \Theta_{\cA}^{\leq r}$. The modular
characteristic, the $r$-matrix, the shadow connection, and the full
genus tower are projections of $\Theta_{\cA}$.


\subsection{Shadow connection}\label{subsec:shadow-connection}

\begin{definition}[Shadow connection]\label{def:shadow-connection}
The \emph{modular shadow connection} on $\overline{\cM}_{g,n}$ is
\begin{equation}\label{eq:shadow-connection}
\nabla^{\hol}_{g,n}
\;=\;
d - \Sh_{g,n}(\Theta_{\cA}),
\end{equation}
where $\Sh_{g,n}$ is the shadow projection to the
$(g,n)$-component of the modular convolution algebra.
\end{definition}

Flatness of $\nabla^{\hol}$ is the MC equation projected to
$(g,n)$. At genus $0$, $\nabla^{\hol}_{0,n}$ is
unconditionally flat; at genus $g \geq 1$, the connection acquires
curvature $\kappa(\cA) \cdot \omega_g$ (the genus-$g$ anomaly).

The shadow connection encodes the full genus tower. Its
monodromy representations are the physical amplitudes: at genus $0$,
the KZ monodromy for affine algebras; the BPZ monodromy for
Virasoro.


\subsection{Collision filtration and shadow depth}
\label{subsec:collision}

The modular convolution algebra $\gmod$ carries a collision
filtration $F^0_{\coll} \supset F^1_{\coll} \supset \cdots$
indexed by minimum pairwise distance on $\FM_n(\bC)$. Depth $0$
is free propagation; depth $k$ is $k$-fold nested collision. The
spectral sequence
$E_1^{p,*} = H^*(\gr^p_{\coll}\gmod, d_0)$ converges to
$H^*(\gmod)$. Collapse is governed by shadow depth:
\begin{center}
\renewcommand{\arraystretch}{1.15}
\begin{tabular}{lccl}
\toprule
Class & Shadow depth & OPE pole order & Representative \\
\midrule
$\mathbf{G}$ & $2$ & $2$ & Heisenberg $\cH_k$ \\
$\mathbf{L}$ & $3$ & $3$ & Affine $\widehat{\fg}_k$ \\
$\mathbf{C}$ & $4$ & $4$ & $\beta\gamma$ system \\
$\mathbf{M}$ & $\infty$ & $\geq 4$ & Virasoro, $\cW_N$ \\
\bottomrule
\end{tabular}
\end{center}
The discriminant $\Delta = 8\kappa S_4$ controls truncation:
$\Delta = 0$ (classes $\mathbf{G}$, $\mathbf{L}$, $\mathbf{C}$)
gives a finite shadow tower; $\Delta \neq 0$ (class $\mathbf{M}$)
gives an infinite tower. The discriminant is linear in $\kappa$,
not quadratic.


% ====================================================================
\section{The Heisenberg datum}\label{sec:heisenberg}
% ====================================================================

The Heisenberg algebra $\cH_k$ at level $k$ is the abelian current
algebra generated by a single field $J$ with OPE
\begin{equation}\label{eq:heis-ope}
J(z)\,J(w) \;\sim\; \frac{k}{(z-w)^2}.
\end{equation}
This is abelian Chern--Simons for $G = U(1)$ at level $k$.

\begin{proposition}[Heisenberg datum]\label{prop:heis-datum}
The holographic datum of $\cH_k$ is
\begin{equation}\label{eq:heis-datum}
\cH(\cH_k) \;=\;
\bigl(\cH_k,\; \Sym^{\mathrm{ch}}(V^*),\;
\mathsf{Vect},\; k/z,\; k \cdot \eta \otimes \Lambda,\;
d\bigr).
\end{equation}
\end{proposition}

\begin{proof}
We verify each component.

\smallskip\noindent\emph{(i) Boundary.}\enspace
$\cA = \cH_k$, the free boson at level $k$.

\smallskip\noindent\emph{(ii) Koszul dual.}\enspace
$\cA^! = \Sym^{\mathrm{ch}}(V^*)$. This is the
curved symmetric chiral algebra with $m_0 = -k\omega$; it is
\emph{not} $\cH_{-k}$ (the same algebra has the same $\kappa$
on both sides, but the duals differ as curved objects).

\smallskip\noindent\emph{(iii) Lines.}\enspace
$\cC = \mathsf{Vect}$. The line category is trivial: abelian
gauge theory has no nontrivial Wilson lines.

\smallskip\noindent\emph{(iv) $r$-matrix.}\enspace
% PE-1 r-matrix check:
% family: Heis
% r(z) written: k/z
% level parameter: k
% OPE pole order: 2
% r-matrix pole order: 1 (d-log absorbs one pole)
% AP126 check: r(z)|_{k=0} = 0. Match: Y.
$r(z) = k/z$. The OPE has a double pole; $d\log$ absorbs one,
leaving a simple pole. At $k = 0$: $r = 0$ (abelian limit).

\smallskip\noindent\emph{(v) MC element.}\enspace
$\Theta_{\cH_k} = k \cdot \eta \otimes \Lambda$, where
$\eta = d\log(z_1 - z_2)$ is the Arnold form and $\Lambda$ is the
unique (up to scale) bilinear on $V$. The MC equation is
satisfied because $\cH_k$ is class $\mathbf{G}$: all operations
beyond degree $2$ vanish.

\smallskip\noindent\emph{(vi) Shadow connection.}\enspace
$\nabla^{\hol} = d$. The shadow connection is flat and trivial:
the MC element is scalar, so $\Sh_{g,n}(\Theta_{\cH_k})$ is
a multiple of the identity for all $(g,n)$.
\end{proof}

\begin{remark}[Modular characteristic]\label{rem:heis-kappa}
% PE-2 kappa check:
% family: Heis
% kappa = k
% at k=0: kappa=0
The modular characteristic is $\kappa(\cH_k) = k$. This is the
abelian prototype: $\av(r(z)) = \av(k/z) = k = \kappa$.
The Koszul complementarity is $\kappa(\cH_k) + \kappa(\cH_k^!) = 0$
(class $\mathbf{G}$/$\mathbf{L}$/free-field complementarity).
\end{remark}

\begin{remark}[Quantised $R$-matrix]\label{rem:heis-R}
The quantised $R$-matrix is $R(z) = \exp(k\hbar/z)$, scalar and
automatically satisfying the Yang--Baxter equation. The monodromy
$\exp(-2\pi i k)$ is the Aharonov--Bohm phase.
\end{remark}


% ====================================================================
\section{The Kac--Moody datum}\label{sec:km}
% ====================================================================

The affine Kac--Moody algebra $\widehat{\fg}_k$ at level $k$
is generated by currents $J^a(z)$ ($a = 1, \ldots, \dim(\fg)$)
with OPE
\begin{equation}\label{eq:km-ope}
J^a(z)\,J^b(w)
\;\sim\;
\frac{k\,\kappa^{ab}}{(z-w)^2}
\;+\;
\frac{f^{ab}{}_c\,J^c(w)}{z-w},
\end{equation}
where $\kappa^{ab}$ is the Killing form and $f^{ab}{}_c$ are
structure constants. This is 3d Chern--Simons for $G$ at level $k$.
We take $\fg = \fsl_2$ for concreteness; the formulae generalise to
arbitrary simple $\fg$.

\begin{proposition}[Affine $\widehat{\fsl}_2$ datum]
\label{prop:km-datum}
The holographic datum of $\widehat{\fsl}_2{}_k$ is
\begin{equation}\label{eq:km-datum}
\cH(\widehat{\fsl}_2{}_k) \;=\;
\Bigl(
\widehat{\fsl}_2{}_k,\;
\widehat{\fsl}_2{}_{-k-4},\;
\cO_q^{\mathrm{sh}},\;
\frac{k\,\Omega}{z},\;
\Theta_k,\;
d - \kappa_k\,\omega_g - \cdots
\Bigr),
\end{equation}
with $q = e^{i\pi/(k+2)}$ and
$\kappa_k = \dim(\fsl_2)(k+h^\vee)/(2h^\vee) = 3(k+2)/4$.
\end{proposition}

\begin{proof}
We verify each component.

\smallskip\noindent\emph{(i) Boundary.}\enspace
$\cA = \widehat{\fsl}_2{}_k$, the affine Kac--Moody vertex
algebra at level $k$. The three generators $J^a$
($a \in \{e, f, h\}$) produce a triple-pole OPE (double pole from
the Killing form, simple pole from the Lie bracket).

\smallskip\noindent\emph{(ii) Koszul dual.}\enspace
$\cA^! = \widehat{\fsl}_2{}_{-k-4}$. The Koszul involution acts by
$k \mapsto -k - 2h^\vee$ with $h^\vee(\fsl_2) = 2$. The
complementarity sum is
$\kappa(\widehat{\fsl}_2{}_k) +
\kappa(\widehat{\fsl}_2{}_{-k-4}) = 0$,
verified at $k = 0$: $\kappa_0 = 3/2$ and $\kappa_{-4} = -3/2$.

\smallskip\noindent\emph{(iii) Lines.}\enspace
$\cC = \cO_q^{\mathrm{sh}} \simeq
\Rep(U_q(\fsl_2))$. The braided tensor category of line
operators is the representation category of the quantum group at
$q = e^{i\pi/(k+2)}$. The identification is the affine
Drinfeld--Kohno theorem: the KZ monodromy representation on
conformal blocks agrees with the quantum-group $R$-matrix action
on tensor products of evaluation modules.

\smallskip\noindent\emph{(iv) $r$-matrix.}\enspace
% PE-1 r-matrix check:
% family: affine KM
% r(z) written: k*Omega/z
% level parameter: k
% convention: trace-form
% OPE pole order: 3 (double from Killing, simple from bracket)
% r-matrix pole order: 2 ... but wait: the r-matrix is k*Omega/z (simple pole)
% because we take the degree-2 collision residue which picks out
% the leading OPE structure after dlog absorption
% AP126 check: r(z)|_{k=0} = 0. Match: Y.
% critical level: r(z)|_{k=-2} = -2*Omega/z (finite, trace-form).
$r(z) = k\,\Omega_{\mathrm{tr}}/z$ in trace-form convention,
where $\Omega_{\mathrm{tr}}$ is the trace-form Casimir. At $k = 0$:
$r = 0$ (abelian limit). In Kazhdan--Lusztig (KZ) normalisation:
$r(z) = \Omega_{\mathrm{KZ}}/((k + h^\vee)z)$ with
$\Omega_{\mathrm{KZ}} = 2 h^\vee\, \Omega_{\mathrm{tr}}$ the
Killing-normalised Casimir. The two conventions are related by a
\emph{level-dependent rescaling of the Casimir}, not by the naive
identity $k\,\Omega_{\mathrm{tr}} = \Omega_{\mathrm{KZ}}/(k+h^\vee)$:
the latter fails as a rational-function equation in $k$ (at $k=0$
the left side vanishes, the right side is finite and non-zero).
The correct statement is that $r_{\mathrm{tr}}$ and $r_{\mathrm{KZ}}$
are two presentations of the same underlying operadic datum in
different bases, exchanged by $\Omega \mapsto \Omega$ and
$k \mapsto k/((k+h^\vee)\cdot 2 h^\vee)$. At $k = -h^\vee = -2$:
trace-form finite, KZ diverges (Sugawara singularity at critical level).

\smallskip\noindent\emph{(v) MC element.}\enspace
$\Theta_k = \Theta_{\widehat{\fsl}_2{}_k}$ is determined by
$D_{\cA}^2 = 0$. The shadow tower terminates at depth $3$
(class $\mathbf{L}$): the quartic invariant $S_4 = 0$
(the discriminant $\Delta = 8\kappa S_4 = 0$), so the shadow
obstruction tower is finite.

\smallskip\noindent\emph{(vi) Shadow connection.}\enspace
At genus $0$, degree $n$, the shadow connection is the KZ
connection:
\begin{equation}\label{eq:kz}
\nabla^{\hol}_{0,n}
\;=\;
d - \sum_{i < j} r_{ij}\,dz_{ij}
\;=\;
d - \frac{1}{k + h^\vee}\sum_{i < j}
\frac{\Omega_{ij}\,dz_{ij}}{z_{ij}},
\end{equation}
where $dz_{ij} = d(z_i - z_j)$. Flatness of the KZ connection
is the MC equation at genus $0$, degree $n$. At genus $g \geq 1$,
the connection acquires curvature
$\kappa_k \cdot \omega_g$ with
$\kappa_k = 3(k+2)/4$.
\end{proof}

\begin{remark}[Quantum group line category]\label{rem:km-qg}
The identification $\cC \simeq \Rep(U_q(\fsl_2))$ is a
reconstruction theorem: the monodromy of the genus-$0$ shadow
connection $\nabla^{\hol}_{0,n}$ (the KZ connection) generates a
braided tensor category that is equivalent to the representation
category of the quantum group. The spectral parameter $z$ in
$r(z) = k\,\Omega/z$ is the residual parameter after
$d\log$-absorption from the OPE; on evaluation modules, it becomes
the additive spectral parameter of the Yangian. The full quantised
$R$-matrix $R(z) = 1 + \hbar\,r(z) + O(\hbar^2)$ satisfies the
quantum Yang--Baxter equation.
\end{remark}

\begin{remark}[Critical level]\label{rem:critical}
At $k = -h^\vee = -2$: $\kappa = 0$, all monodromy is trivial,
$H^1$ doubles (from $4$ to $8$ for $\fsl_2$), and Koszulness
fails. The Feigin--Frenkel centre is the algebra of
$\fsl_2$-opers on the formal disk: $Z(V_{-2}(\fsl_2)) \simeq
\mathrm{Fun}(\mathrm{Op}_{\fsl_2}(D))$. The shadow connection
degenerates, and the line category jumps from semisimple to
abelian.
\end{remark}


% ====================================================================
\section{The Virasoro datum}\label{sec:virasoro}
% ====================================================================

The Virasoro algebra $\Vir_c$ at central charge $c$ is generated by
the stress tensor $T$ with OPE
\begin{equation}\label{eq:vir-ope}
T(z)\,T(w)
\;\sim\;
\frac{c/2}{(z-w)^4}
\;+\;
\frac{2T(w)}{(z-w)^2}
\;+\;
\frac{\partial T(w)}{z-w}.
\end{equation}
This is 3d gravity: $SL(2,\bR) \times SL(2,\bR)$ Chern--Simons
at levels $k_\pm = \ell/(4G)$, with $c = 3\ell/(2G)$
(Brown--Henneaux).

\begin{proposition}[Virasoro datum]\label{prop:vir-datum}
The holographic datum of $\Vir_c$ is
\begin{equation}\label{eq:vir-datum}
\cH(\Vir_c)
\;=\;
\bigl(\Vir_c,\; \Vir_{26-c},\; \cC_c,\; r_c(z),\;
\Theta_c,\; \nabla^{\hol}_c\bigr),
\end{equation}
with
\begin{align}
r_c(z) &\;=\; \frac{c/2}{z^3} + \frac{2T}{z},
\label{eq:vir-r}\\[3pt]
\kappa_c &\;=\; \frac{c}{2},
\label{eq:vir-kappa}\\[3pt]
\nabla^{\hol}_{1,1}
&\;=\; d \;-\; \kappa_c\,\omega_1
\;-\; \frac{10}{c(5c+22)}\,\delta_4 \;-\; \cdots.
\label{eq:vir-connection}
\end{align}
\end{proposition}

\begin{proof}
We verify each component.

\smallskip\noindent\emph{(i) Boundary.}\enspace
$\cA = \Vir_c$, the Virasoro vertex algebra.

\smallskip\noindent\emph{(ii) Koszul dual.}\enspace
$\cA^! = \Vir_{26-c}$. The Koszul involution is
$c \mapsto 26 - c$. The complementarity sum is
$\kappa(\Vir_c) + \kappa(\Vir_{26-c}) = c/2 + (26-c)/2 = 13$.
Self-dual at $c = 13$ (not $c = 26$).

\smallskip\noindent\emph{(iii) Lines.}\enspace
$\cC_c$ is the line-operator category, determined by the
$r$-matrix monodromy. For the Virasoro algebra, the line
category is infinite-depth: the monodromy representation is
determined by the full BPZ differential equations, not by a
finite-dimensional quantum group.

\smallskip\noindent\emph{(iv) $r$-matrix.}\enspace
% PE-1 r-matrix check:
% family: Vir
% r(z) written: (c/2)/z^3 + 2T/z
% level parameter: c
% OPE pole order: 4
% r-matrix pole order: 3 (d-log absorbs one)
% Convention: the r-matrix has cubic + simple, NOT quartic
The $r$-matrix $r_c(z) = (c/2)/z^3 + 2T/z$ has a cubic pole
(gravitational, from the quartic OPE pole after $d\log$-absorption)
and a simple pole (stress-tensor exchange). The quartic contact
invariant
\begin{equation}\label{eq:quartic}
Q^{\mathrm{contact}}
\;=\;
\frac{10}{c(5c+22)}
\;\neq\; 0
\end{equation}
is nonzero for all $c \neq 0$, $c \neq -22/5$: this is the
obstruction to extending $r(z)$ from degree $2$ to degree $4$, and
its non-vanishing is the signature of class $\mathbf{M}$.

\smallskip\noindent\emph{(v) MC element.}\enspace
$\Theta_c = \Theta_{\Vir_c}$ is the infinite-depth MC element.
The shadow tower does not truncate: all shadow invariants
$S_r$ ($r \geq 2$) are nonzero.
% PE-2 kappa check:
% family: Vir
% kappa = c/2
% at c=0: kappa=0
% at c=13: kappa=13/2 (self-dual)
The modular characteristic is
$\kappa(\Vir_c) = S_2 = c/2$.

\smallskip\noindent\emph{(vi) Shadow connection.}\enspace
The shadow connection~\eqref{eq:vir-connection} has leading
curvature $\kappa_c\,\omega_g = (c/2)\,\omega_g$ at genus $g$,
plus an infinite series of corrections from the shadow
obstruction tower.
\end{proof}

\begin{remark}[The gravitational hierarchy]\label{rem:vir-gravity}
Shadow class $\mathbf{M}$ means the $\Ainf$ tower
$(m_1, m_2, m_3, \ldots)$ is infinite: every higher operation
$m_k$ is nonzero. The generating depth is $d_{\mathrm{gen}} = 3$
(the operation $m_3$ determines all $m_k$ recursively via the MC
equation), but the algebraic depth is
$d_{\mathrm{alg}} = \infty$ (no operation vanishes). Gravity
does not truncate.
\end{remark}

\begin{remark}[Two critical values of $c$]\label{rem:vir-critical}
The two critical central charges play distinct roles:
\begin{enumerate}[nosep,label=(\roman*)]
\item $c = 13$: Koszul self-duality
  ($\Vir_{13}^! = \Vir_{13}$). The complementarity parameter is
  $\kappa + \kappa' = 13 \neq 0$; self-duality is a symmetry
  enhancement within the perturbative regime.
\item $c = 26$: anomaly cancellation
  ($\kappa(\Vir_{26-c}^!) = (26-c)/2 = 0$ at $c = 26$). The
  Clifford/exterior bifurcation
  (Section~\ref{sec:anomaly}). This is the bosonic string.
\end{enumerate}
Self-duality and anomaly cancellation are separated by $13$
units of central charge.
\end{remark}

\begin{remark}[Primitive coproduct]\label{rem:vir-coproduct}
The transferred coproduct $\Delta_z^{\Vir}$ is strictly primitive
at all degrees: the ghost-number obstruction intrinsic to
Drinfeld--Sokolov reduction kills every higher-product-lattice
tree correction. All gravitational dynamics resides in the
collision residue and the infinite $\Ainf$ product tower; the
coproduct channel is identically zero.
\end{remark}


% ====================================================================
\section{Comparison table}\label{sec:comparison}
% ====================================================================

\begin{center}
\renewcommand{\arraystretch}{1.3}
\begin{tabular}{@{}lccc@{}}
\toprule
Component & $\cH_k$ & $\widehat{\fsl}_2{}_k$
 & $\Vir_c$ \\
\midrule
$\cA$ (boundary)
 & $\cH_k$
 & $\widehat{\fsl}_2{}_k$
 & $\Vir_c$ \\[3pt]
$\cA^!$ (Koszul dual)
 & $\Sym^{\mathrm{ch}}(V^*)$
 & $\widehat{\fsl}_2{}_{-k-4}$
 & $\Vir_{26-c}$ \\[3pt]
$\cC$ (lines)
 & $\mathsf{Vect}$
 & $\Rep(U_q(\fsl_2))$
 & $\cC_c$ (BPZ) \\[3pt]
$r(z)$
 & $k/z$
 & $k\,\Omega/z$
 & $(c/2)/z^3 + 2T/z$ \\[3pt]
$\kappa$
 & $k$
 & $3(k+2)/4$
 & $c/2$ \\[3pt]
$\kappa + \kappa'$
 & $0$
 & $0$
 & $13$ \\[3pt]
Class
 & $\mathbf{G}$
 & $\mathbf{L}$
 & $\mathbf{M}$ \\[3pt]
Depth
 & $2$ & $3$ & $\infty$ \\[3pt]
$Q^{\mathrm{contact}}$
 & $0$ & $0$
 & $10/[c(5c+22)]$ \\[3pt]
$\nabla^{\hol}$
 & $d$ (trivial) & KZ & BPZ (infinite) \\[3pt]
Physics
 & abelian CS & non-abelian CS & 3d gravity \\
\bottomrule
\end{tabular}
\end{center}

The three families represent three qualitatively distinct regimes.
Class $\mathbf{G}$ is Gaussian: scalar, trivial lines, the datum
collapses to a single number $k$. Class $\mathbf{L}$ is Lie-theoretic:
matrix-valued, finite quantum-group lines, the datum is determined by
$\fg$ and $k$. Class $\mathbf{M}$ is gravitational: infinite depth,
no finite quantum group, the datum requires the full infinite tower
$\Theta_c$.


% ====================================================================
\section{The BBL triangle}\label{sec:bbl}
% ====================================================================

The bulk--boundary--line (BBL) triangle is the central diagram of 3d
holomorphic-topological gauge theory. It is not built from three
independent objects connected by ad hoc functors; it is a single
object (the bar complex with its $\Eone$-chiral coassociative
structure) viewed from three angles.


\subsection{The three vertices}\label{subsec:vertices}

\begin{definition}[BBL triangle]\label{def:bbl}
Let $\cA$ be a chirally Koszul vertex algebra. The
\emph{bulk--boundary--line triangle} of $\cA$ consists of:
\begin{enumerate}[label=(\roman*),nosep]
\item \emph{Boundary}: the chiral algebra $\cA$ itself, the
  input to the bar construction.
\item \emph{Bulk}: the chiral derived centre
  \begin{equation}\label{eq:bulk}
  \cA_{\mathrm{bulk}}
  \;\simeq\;
  \cZ^{\mathrm{der}}_{\mathrm{ch}}(\cA)
  \;=\;
  C^\bullet_{\mathrm{ch}}(\cA, \cA),
  \end{equation}
  computed by chiral Hochschild cochains (not topological Hochschild;
  the geometry of the curve determines which
  chiral/topological/categorical Hochschild theory applies).
\item \emph{Lines}: the category
  $\cC \simeq \cA^!\text{-}\mathsf{mod}$ of modules over the
  chiral Koszul dual.
\end{enumerate}
\end{definition}

The $\SCchtop$ structure lives on the \emph{pair}
$(\cZ^{\mathrm{der}}_{\mathrm{ch}}(\cA), \cA)$: the chiral
derived centre carries brace operations and a chiral Gerstenhaber
bracket, and the boundary algebra $\cA$ is acted upon by the bulk
via the universal brace. The SC datum is two-coloured with
directionality: bulk-to-boundary only, no boundary-to-bulk.


\subsection{The three edges}\label{subsec:edges}

\begin{enumerate}[label=(\roman*)]
\item \emph{Boundary $\to$ Bulk} (the bar map).
The bar construction maps $\cA$ to $\barB(\cA)$; the trace map
$\tr_g \colon \barB^{(g)}(\cA) \to \mathbf{C}_g(\cA)$ projects
the bar coalgebra into the ambient complex
$\mathbf{C}_g(\cA) = R\Gamma(\overline{\cM}_g,
\cZ(\cA))$. The image
$\mathbf{Q}_g(\cA) = \mathrm{im}(\tr_g)$ is one half of the
complementarity decomposition.

\item \emph{Bulk $\to$ Lines} (the complementarity cofibre).
The complementarity theorem gives
\[
\mathbf{C}_g(\cA) \;\simeq\;
\mathbf{Q}_g(\cA) \;\oplus\; \mathbf{Q}_g(\cA^!).
\]
The cofibre $\cofib(\mathbf{Q}_g(\cA) \hookrightarrow
\mathbf{C}_g(\cA)) \simeq \mathbf{Q}_g(\cA^!)$ is the dual
shadow, controlling the line-operator category through the
centre--module adjunction.

\item \emph{Boundary $\to$ Lines} (Koszul duality).
The composite $\cA \mapsto \barB(\cA) \mapsto H^*(\barB(\cA))
\mapsto (H^*(\barB(\cA)))^\vee = \cA^!$ is bar-dualize. The
module category $\cA^!\text{-}\mathsf{mod}$ is the categorical
output.
\end{enumerate}

Commutativity of the triangle is Theorem~A (the Verdier
intertwining at the algebra level).


\subsection{Reconstruction from the datum}\label{subsec:reconstruction}

The datum $\cH(\cA)$ is a complete invariant.

\begin{proposition}[Recovery theorems]\label{prop:recovery}
The following reconstructions hold:
\begin{enumerate}[label=\textup{(R\arabic*)},nosep]
\item \emph{Collision residue is twisting.}
  $\Res^{\coll}_{0,2}(\Theta_{\cA}) = \tau|_{\barB^2(\cA)}$
  (the universal twisting morphism restricted to bar degree $2$).
\item \emph{CYBE from Arnold.}
  The three boundary divisors of $\overline{M}_{0,4}$ produce the
  three terms of $\sum [r_{ij}, r_{ik}] = 0$.
\item \emph{Affine shadow = KZ.}
  For $\cA = \widehat{\fg}_k$:
  $\Sh_{0,n}(\Theta_{\widehat{\fg}_k})
  = (k + h^\vee)^{-1} \sum_{i<j}
  \Omega_{ij}\,dz_{ij}/z_{ij}$.
\item \emph{Quartic as $L_\infty$ obstruction.}
  $[\mathfrak{R}^{\mathrm{mod}}_4(\cA)]
  = [o_4(\cA)] \in
  H^2(\cA^{\mathrm{sh}}_{4,0})$ is the obstruction to extending
  $r(z)$ from degree $2$ to degree $4$. Vanishes for
  $\mathbf{G}$, $\mathbf{L}$, $\mathbf{C}$; nonzero for
  $\mathbf{M}$.
\end{enumerate}
\end{proposition}

\begin{remark}[BBL for Heisenberg]\label{rem:bbl-heis}
Boundary $= \cH_k$. Bulk
$= \cZ^{\mathrm{der}}_{\mathrm{ch}}(\cH_k)$, which is
$3$-dimensional (the naive centre is $1$-dimensional; the derived
centre sees $\mathrm{Ext}^1$ and $\mathrm{Ext}^2$). Lines
$= \mathsf{Vect}$. The triangle degenerates: the bulk-to-lines edge
is trivial, and the boundary-to-lines edge is trivial.
\end{remark}

\begin{remark}[BBL for Kac--Moody]\label{rem:bbl-km}
Boundary $= \widehat{\fsl}_2{}_k$. Bulk
$= \cZ^{\mathrm{der}}_{\mathrm{ch}}(\widehat{\fsl}_2{}_k)$, with
$\ChirHoch^1(V_k(\fsl_2)) \cong \fsl_2$ and total dimension
$\dim(\fsl_2) + 2 = 5$. Lines $= \Rep(U_q(\fsl_2))$. The
Drinfeld--Kohno theorem makes the boundary-to-lines edge
constructive: the monodromy of the KZ connection generates
$\Rep(U_q(\fsl_2))$. The bulk-to-lines edge is the complementarity
cofibre: $\mathbf{Q}_g(\widehat{\fsl}_2{}_{-k-4})$ controls the
line-operator anomaly.
\end{remark}

\begin{remark}[BBL for Virasoro]\label{rem:bbl-vir}
Boundary $= \Vir_c$. Bulk
$= \cZ^{\mathrm{der}}_{\mathrm{ch}}(\Vir_c)$: chiral Hochschild
cohomology concentrated in degrees $\{0, 1, 2\}$ with polynomial
Hilbert series (this is amplitude, not virtual dimension). Lines
$= \cC_c$: the infinite-depth line category, governed by the BPZ
equations. The triangle is fully active at all three vertices;
the bulk-to-lines edge is controlled by the Virasoro
complementarity $\mathbf{Q}_g(\Vir_c) \oplus
\mathbf{Q}_g(\Vir_{26-c})$.
\end{remark}


% ====================================================================
\section{The anomaly completion}\label{sec:anomaly}
% ====================================================================


\subsection{The transgression algebra}\label{subsec:transgression}

Fix a chirally Koszul algebra $\cA$ with Koszul dual
$B = \cA^!$. The curvature of the bar complex on a genus-$g$
curve is $d^2_{\mathrm{fib}} = \kappa(\cA) \cdot \omega_g$. The
Koszul dual carries dual curvature
$\lambda = \kappa(B) \cdot \omega_g$.

\begin{definition}[Transgression algebra]\label{def:transgression}
The \emph{transgression algebra} is
\begin{equation}\label{eq:transgression}
B_\lambda
\;=\;
B\langle\eta\rangle \,\big/\,
(\eta^2 - \lambda),
\qquad
|\eta| = 1,
\quad
d_{B_\lambda}(\eta) = 0,
\end{equation}
where $\eta$ is a degree-$1$ generator with Clifford (not
exterior) relation $\eta^2 = \lambda = \kappa(B) \cdot \omega_g$,
and the differential on $B$ is extended by
$d_{B_\lambda}(b) = d_B(b) + [\eta, b]$ for $b \in B$.
\end{definition}

The relation $\eta^2 = \lambda$ is Clifford when $\lambda \neq 0$
and exterior when $\lambda = 0$. This is the algebraic encoding of
the anomaly: $\lambda = 0$ means the theory is anomaly-free, and
the algebra structure changes qualitatively.


\subsection{The secondary anomaly}\label{subsec:secondary}

\begin{definition}[Secondary anomaly]\label{def:secondary}
The \emph{secondary anomaly} is
\begin{equation}\label{eq:secondary}
u
\;=\;
\eta^2
\;=\;
\lambda
\;=\;
\kappa(\cA^!) \cdot \omega_g.
\end{equation}
It is closed, central, and acts as the scalar $\kappa(\cA^!)$
on standard modules.
\end{definition}

The secondary anomaly is the bifurcation parameter for the genus
tower. Its behaviour depends on the Koszul complementarity:

\begin{center}
\renewcommand{\arraystretch}{1.15}
\begin{tabular}{@{}lccl@{}}
\toprule
Family & $\kappa(\cA^!)$ & $u$ & Regime \\
\midrule
$\cH_k$ & $-k$ & $-k\,\omega_g$ & $u = 0$ at $k = 0$ \\
$\widehat{\fsl}_2{}_k$ & $3(-k-2)/4$ & $3(-k-2)\omega_g/4$ &
  $u = 0$ at $k = -2$ (critical) \\
$\Vir_c$ & $(26-c)/2$ & $(26-c)\omega_g/2$ &
  $u = 0$ at $c = 26$ (string) \\
\bottomrule
\end{tabular}
\end{center}


\subsection{Genus-$g$ Clifford completion}\label{subsec:clifford}

\begin{construction}[Clifford completion]\label{con:clifford}
At genus $g$, adjoin $2g$ handle operators
$\alpha_1, \beta_1, \ldots, \alpha_g, \beta_g$ with
\begin{equation}\label{eq:clifford-rel}
\alpha_i \beta_j + \beta_j \alpha_i = \delta_{ij}\,u,
\qquad
\alpha_i\alpha_j + \alpha_j\alpha_i = 0,
\qquad
\beta_i\beta_j + \beta_j\beta_i = 0.
\end{equation}
The genus-$g$ Clifford completion is the algebra
\begin{equation}\label{eq:clifford-completion}
\mathfrak{G}_g(\cA)
\;=\;
B_\lambda \otimes \mathrm{Cl}(H_1(\Sigma_g); u).
\end{equation}
\end{construction}

The handle operators record the genus-$g$ topology. Their
anticommutator is controlled by the secondary anomaly $u$, which
splits quantum field theory into two qualitative regimes.


\subsection{The dichotomy}\label{subsec:dichotomy}

\begin{proposition}[Clifford/exterior dichotomy]
\label{prop:dichotomy}
The genus tower of $\cA$ bifurcates at $u = 0$:
\begin{enumerate}[label=(\roman*),nosep]
\item When $u \neq 0$: the Clifford algebra
  $\mathrm{Cl}(H_1(\Sigma_g); u)$ is a matrix algebra (dimension
  $2^g$) after inverting $u$. Genus-$g$ is Morita equivalent to
  genus-$0$. All genus corrections are resummable from genus-$0$
  data. This is the perturbative regime.
\item When $u = 0$: the Clifford anticommutator vanishes. The
  handle operators become exterior generators on $H_1(\Sigma_g)$
  (dimension $2^{2g}$). The theory sees genuine surface topology.
  The genus tower collapses ($\kappa_{\mathrm{eff}} = 0$).
\end{enumerate}
\end{proposition}

\begin{proof}
Standard Clifford algebra theory. When the bilinear form
$u \cdot \langle\cdot,\cdot\rangle$ on $H_1(\Sigma_g, \bZ)$ is
nondegenerate (i.e., $u \neq 0$), the Clifford algebra is a
matrix algebra $M_{2^g}(\bC)$ (even-dimensional, nondegenerate
symplectic). When $u = 0$, the bilinear form degenerates and the
Clifford algebra becomes the exterior algebra
$\Lambda^*(H_1(\Sigma_g)^*)$.
\end{proof}


\subsection{The three families}\label{subsec:anomaly-families}

\noindent\emph{Heisenberg.}\enspace
$u = -k\,\omega_g$. At $k = 0$: $u = 0$, exterior regime. At
generic $k$: $u \neq 0$, perturbative. The bifurcation is at the
abelian limit.

\smallskip
\noindent\emph{Kac--Moody.}\enspace
$u = 3(-k-2)\omega_g/4$. At $k = -2$ (critical level): $u = 0$,
the Feigin--Frenkel centre appears, the genus tower collapses. At
generic $k$: $u \neq 0$, perturbative. The critical level is
the bifurcation point.

\smallskip
\noindent\emph{Virasoro.}\enspace
$u = (26-c)\omega_g/2$. At $c = 26$: $u = 0$, the bosonic string
becomes topological ($Z_{\mathrm{grav}} = 1$). At generic $c$:
$u \neq 0$, perturbative. At $c = 13$: $u = 13\omega_g/2 \neq 0$,
so Koszul self-duality lives within the perturbative regime.

The gravitational partition function at genus $g$ is the
super-trace of the MC element through the Clifford factor:
\begin{equation}\label{eq:partition}
Z_g
\;=\;
\mathrm{str}_{\mathfrak{G}_g / B_\Theta}(\Theta^{(g)}_c).
\end{equation}
The complementarity theorem gives a Lagrangian decomposition
\begin{equation}\label{eq:lagrangian}
H^*(\overline{\cM}_g,\, \cZ(\Vir_c))
\;\simeq\;
\mathbf{Q}_g(\Vir_c)
\;\oplus\;
\mathbf{Q}_g(\Vir_{26-c})
\end{equation}
into boundary ($+1$-eigenspace of Verdier involution) and bulk
($-1$-eigenspace) sectors. The shifted symplectic structure makes
both Lagrangian.


\subsection{The anomaly parameter as holographic observable}
\label{subsec:anomaly-observable}

The secondary anomaly $u$ is the single scalar that governs the
passage from genus $0$ to all genera. For each family, $u$ vanishes
at a distinguished value of the level/central charge:

\begin{center}
\renewcommand{\arraystretch}{1.15}
\begin{tabular}{@{}lccl@{}}
\toprule
Family & $u = 0$ at & Physical meaning & Regime \\
\midrule
$\cH_k$ & $k = 0$ & trivial theory & exterior \\
$\widehat{\fg}_k$ & $k = -h^\vee$ & critical level, Feigin--Frenkel &
  exterior \\
$\Vir_c$ & $c = 26$ & bosonic string, anomaly cancellation &
  exterior \\
\bottomrule
\end{tabular}
\end{center}

\noindent
The datum $\cH(\cA)$ together with the anomaly completion
$B_\Theta$ controls the full 3d HT correspondence: the six
components determine the genus-$0$ physics (the datum), and the
secondary anomaly $u$ determines the genus tower (the completion).
The entire structure descends from a single identity: $D_{\cA}^2 = 0$.


% ====================================================================
\section{Concluding remarks}\label{sec:conclusion}
% ====================================================================

\subsection{The $E_n$ circle}\label{subsec:en-circle}

The holographic datum sits inside the $\En$ operadic circle:
\begin{multline*}
\Ethree^{\mathrm{top}}(\text{bulk})
\;\to\;
\Etwo(\text{boundary chiral})
\;\to\;
\Eone(\text{bar/QG})\\
\;\to\;
\Etwo(\text{Drinfeld centre})
\;\to\;
\Ethree^{\mathrm{top}}(\text{derived centre}).
\end{multline*}
Each arrow is: restriction to codimension-$2$ defect, ordered bar
complex, categorified averaging (Drinfeld centre), higher Deligne
(derived centre). The circle partly closes: for simple $\fg$, the
$\Ethree$-topological deformation space of
$\cZ^{\mathrm{der}}_{\mathrm{ch}}(V_k(\fg))$ is
$\hbar\,H^3(\fg)[[\hbar]]$, one-dimensional; the formal disk
restriction recovers the Costello--Francis--Gwilliam Chern--Simons
$\Ethree$-algebra. The passage from $\SCchtop$ to $\Ethree$ requires
a conformal vector at non-critical level (Sugawara), which makes
$\bC$-translations $Q$-exact and topologises the holomorphic
direction. Without conformal vector: stuck at $\SCchtop$.

\subsection{Scope}\label{subsec:scope}

The topologisation $\SCchtop \to \Ethree^{\mathrm{top}}$ is proved
for affine Kac--Moody algebras $V_k(\fg)$ at non-critical level
$k \neq -h^\vee$ (where the Sugawara construction is explicit). For
general chiral algebras with conformal vector (Virasoro, $\cW$-algebras),
the topologisation is cohomological but the chain-level statement is
open. The proof is cohomological: for class $\mathbf{M}$ algebras,
where chain-level data is essential, $\Ethree$ may exist only on
cohomology.

The chiral Hochschild cohomology $\ChirHoch^*(\cA)$ is concentrated
in degrees $\{0, 1, 2\}$ with polynomial Hilbert series
(Theorem~H). This is cohomological amplitude, not virtual dimension.
It is not $\bC[\Theta]$, and it is not the Gelfand--Fuchs cohomology
(which is infinite-dimensional).


\subsection{From the datum to physics}\label{subsec:physics}

The derived chiral centre IS the bulk of a real 3d HT gauge theory.
This is the content of the BBL identification: the bar construction
on $\cA$ produces a coalgebra; the chiral Hochschild cochain complex of
$\cA$ (computed from the bar resolution) produces the bulk
observables; the $\SCchtop$ structure on the pair
$(\cZ^{\mathrm{der}}_{\mathrm{ch}}(\cA), \cA)$ encodes the
open/closed interaction. The datum $\cH(\cA)$ is the smallest
package from which the full physical correspondence reconstructs.

The programme beyond this paper: the global BBL triangle (bulk
determines boundary determines bulk, proved boundary-linear; full
reconstruction open), the chain-level topologisation for class
$\mathbf{M}$, and the closing of the $\En$ circle via the
Drinfeld centre identification
$Z(\mathbf{U}_{\cA}) \simeq
\cZ^{\mathrm{der}}_{\mathrm{ch}}(\cA)$.


% ====================================================================
% References
% ====================================================================

\bibliographystyle{amsalpha}
\begin{thebibliography}{99}

\bibitem{BD04}
A.~Beilinson and V.~Drinfeld,
\emph{Chiral Algebras},
AMS Colloquium Publications~\textbf{51}, 2004.

\bibitem{CG17}
K.~Costello and O.~Gwilliam,
\emph{Factorization Algebras in Quantum Field Theory}, Vols.~1--2,
Cambridge University Press, 2017--2021.

\bibitem{CFG25}
K.~Costello, D.~Francis, and D.~Gwilliam,
An $E_3$-algebra from Chern--Simons theory,
arXiv:2602.12412, 2025.

\bibitem{CG18}
K.~Costello and D.~Gaiotto,
Vertex algebras and 4d $\mathcal{N}=2$ theories,
\emph{JHEP} \textbf{2019}(06):018, 2019.

\bibitem{DK90}
V.~Drinfeld and D.~Kazhdan,
On the structure of quantum groups,
\emph{Funct.\ Anal.\ Appl.}\ \textbf{23}(1989), 115--132.

\bibitem{EV98}
P.~Etingof and A.~Varchenko,
Geometry and classification of solutions of the classical dynamical
Yang--Baxter equation,
\emph{Comm.\ Math.\ Phys.}\ \textbf{192}(1998), 77--120.

\bibitem{FF92}
B.~Feigin and E.~Frenkel,
Affine Kac--Moody algebras at the critical level and Gelfand--Dikii
algebras,
\emph{Int.\ J.\ Mod.\ Phys.}\ A\textbf{7}(Suppl.\ 1A)(1992),
197--215.

\bibitem{GZ26}
D.~Gaiotto and J.~Zeng,
Sphere algebras,
arXiv:2026.xxxxx, 2026.

\bibitem{GK94}
V.~Ginzburg and M.~Kapranov,
Koszul duality for operads,
\emph{Duke Math.\ J.}\ \textbf{76}(1994), 203--272.

\bibitem{KZ26}
A.~Khan and J.~Zeng,
The holomorphic-topological Poisson sigma model,
arXiv:2026.xxxxx, 2026.

\bibitem{KZ84}
V.~Knizhnik and A.~Zamolodchikov,
Current algebra and Wess--Zumino model in two dimensions,
\emph{Nuclear Phys.}\ B\textbf{247}(1984), 83--103.

\bibitem{Liv06}
M.~Livernet,
Koszulity of the category of Swiss-cheese operads,
\emph{J.\ Pure Appl.\ Algebra}\ \textbf{204}(2006), 475--495.

\bibitem{LV12}
J.-L.~Loday and B.~Vallette,
\emph{Algebraic Operads},
Grundlehren~\textbf{346}, Springer, 2012.

\bibitem{Lor26}
R.~Lorgat,
\emph{Modular Koszul Duality},
Vols.~I--III, 2026.

\end{thebibliography}

\end{document}
